\documentclass[10pt]{article}
%------------------------Dimensions--------------------------------------------

\headheight=0pt %1in margins at top and bottom (1 inch is added to this value by default)
\headsep=0pt %Increase to increase white space in between headers and the top of the page
\textheight=9.0in %How tall the text body is allowed to be on each page
\usepackage[T1]{fontenc}
\usepackage{lmodern}
\usepackage{microtype}
\usepackage[colorlinks=true, urlcolor={blue}]{hyperref}
\urlstyle{same}
\usepackage{tcolorbox}
\usepackage[top=.5in, bottom=.5in, left=.8in, right=.8in]{geometry}
\usepackage{enumitem}
\setlist{nosep}

\tcbset{
    frame code={}
    center title,
    left=0pt,
    right=0pt,
    top=0pt,
    bottom=0pt,
    colback=gray!40,
    colframe=white,
    width=\dimexpr\textwidth\relax,
    enlarge left by=0mm,
    boxsep=5pt,
    arc=0pt,outer arc=0pt,
    }

\begin{document}
\pagenumbering{gobble}

%%%%%%%%%%%%%%%%%%%%%%%%%%%%%%%%%%%%
\centerline{{\Huge \sc Dan Kluskiewicz}}  %Makes whatever text you put in parenthesis move to the center
%Prevents the following text from being indented
%This is the same as a return in Latex
\centerline{Index, WA \textbullet \hspace{5pt} (215) 421-9496}
\centerline{\href{mailto:dankluskiewicz@gmail.com}{dankluskiewicz@gmail.com} 
\textbullet \hspace{5pt}  \url{https://github.com/dankluskiewicz}
}

\noindent
%\textbf{Objective:} I would like a field-work intensive job in geophysics.

\vspace{-2mm}
\noindent
\rule{\textwidth}{0.4pt}\\
\vspace{-3.5mm}

\begin{small}

%%%%%%%%%%%%%%%%%%%%%%%%%%%%%%
%%%%%%%%%%%%%%%%%%%%%%%%%%%%%% Summary

\noindent {\Large \bf Summary}
\noindent Applied researcher building Python software packages and ML-driven pipelines for multi-terabyte geospatial and remote-sensing datasets (imagery, LiDAR, and field data), delivering validated inventory products with uncertainty estimates.

\vspace{-2mm}
\noindent
\rule{\textwidth}{0.4pt}\\
\vspace{-3.5mm}

%%%%%%%%%%%%%%%%%%%%%%%%%%%%%%
%%%%%%%%%%%%%%%%%%%%%%%%%%%%%% Skills
\noindent {\Large \bf Skills}

\noindent {\bf Software Development:} package development in Python with an emphasis on supporting large, complex, and reproducible workflows, model training and validation, and performance-minded data processing

\noindent {\bf Geospatial \& Remote Sensing:} large-scale spatial-data processing, aerial imagery, LiDAR, spatial sampling

\noindent {\bf ML \& Statistics:} training, deploying, and validating chains of models for object detection, classification, and regression tasks.

\vspace{-2mm}
\noindent
\rule{\textwidth}{0.4pt}\\
\vspace{-3.5mm}

%%%%%%%%%%%%%%%%%%%%%%%%%%%%%%
%%%%%%%%%%%%%%%%%%%%%%%%%%%%%% Jobs
\noindent {\Large \bf Work Experience}

\noindent
\textbf{Barr Geospatial} --- Moscow, ID (working in Index, WA) \hfill \textbf{September 2018 -- Present}

\noindent
\,\,
\textbf{\textit{Programmer Statistician}} 

\noindent
Lead research and development for tools that predict individual-tree forest inventory from a combination of field data, aerial imagery, and LiDAR.

\begin{itemize}[leftmargin = .35in,     noitemsep]
\item Developed a reproducible end-to-end workflow that organizes multi-terabyte remote-sensing data, incorporates field data to train models that detect and classify trees, evaluates models with validation statistics, and scales inference across large areas.
\item Built Python packages to support and implement this workflow, and to estimate and document uncertainty in our inventories.
\item Collaborated with an expanded team of researchers and engineers to maintain and improve this workflow and its supporting tools.
\item Collaborated with foresters and biometricians to design field-sampling protocols that balance cost, logistics, and statistical precision.
\item Presented inventory products and validation results to clients and at conferences (domestic and international).
\end{itemize}

\vspace{1mm}
\noindent
\textbf{Talus Analytics} --- Boulder, CO \hfill \textbf{June 2016 -- April 2018}

\noindent
\,\,
\textbf{\textit{Quantitative Researcher}} \hfill October 2017 -- April 2018

\noindent
\,\,
\textbf{\textit{Junior Quantitative Researcher}} \hfill June 2016 -- September 2017

\begin{itemize}[leftmargin = .35in, 	noitemsep]
\item Developed predictive models for flooding and wildfire spread
\begin{itemize}[leftmargin = .3in, 	noitemsep]
\item Researched existing methods and theory
\item Proposed a geometric method for probabilistic flood-hazard assessment, and led a collaborative effort to implement it in Python.
\item Implemented a simplified static fire model (from scratch), then extended it to include time-dependent predictions
\end{itemize}
\item Developed and implemented a wildfire accumulation model that delineates geographic areas based on wildfire hazards and predicts probable maximum losses within them
\item Wrote the methods section for a FEMA publication on rapid flood modeling
\item Wrote and extensively edited company reports on flood modeling and resilience investments
\end{itemize}

\vspace{-2mm}
\noindent
\rule{\textwidth}{0.4pt}\\
\vspace{-4mm}

%%%%%%%%%%%%%%%%%%%%%%%%%%%%%%
%%%%%%%%%%%%%%%%%%%%%%%%%%%%%% Education
\noindent {\Large \bf Education}

\begin{tcolorbox}
\noindent \centerline{\large \bf University of Washington \hfill 2015} \\
\textbf{M.S. in Geophysics} \hfill GPA: 3.94 \\
{\normalsize \hfill NSF Graduate Research Fellowship}
\end{tcolorbox}
\vspace{-1.5mm}

\vspace{-1mm}
\begin{tcolorbox}
\noindent \centerline{\large \bf Pennsylvania State University, Schreyer Honors College \hfill 2010}
\textbf{B.S. in Physics and Mathematics}\hfill GPA: 3.99\\
{\normalsize \hfill Highest Distinction, Class rank \textbf{1} out of \textbf{519} (College of Science)}
\end{tcolorbox}
\vspace{-1.5mm}

\vspace{-2mm}
\noindent
\rule{\textwidth}{0.4pt}\\
\vspace{-3.5mm}

%%%%%%%%%%%%%%%%%%%%%%%%%%%%%%
%%%%%%%%%%%%%%%%%%%%%%%%%%%%%% Publications
\noindent {\Large \bf Publications}
\begin{itemize}
\item Kluskiewicz, Dan and others. (2017). Sonic methods for measuring crystal orientation fabric in ice, and results from the West Antarctic ice sheet (WAIS) Divide. \textit{Journal of Glaciology}. 63. 1--15. \href{https://doi.org/10.1017/jog.2017.20}{doi:10.1017/jog.2017.20}.
\item Longenecker, Herbert and others. (2019). A Rapid Flood Risk Assessment Method for Response Operations and Non-subject-matter-expert Community Planning. \textit{Journal of Flood Risk Management}. 13. \href{https://doi.org/10.1111/jfr3.12579}{doi:10.1111/jfr3.12579}.
\end{itemize}

\end{small}
\end{document}